% =============================================================================
% FIRE 2026 — Regular Paper (DOUBLE-BLIND)
% ACM ICPS, sigconf, two columns. Maximum 9 pages of content, references excluded.
%
% Anonymity: no author names, no affiliation, no repository URL that identifies
% the authors. Concrete artefact locations are restored in the camera-ready.
% =============================================================================
\documentclass[sigconf,natbib=true,anonymous=true]{acmart}

\AtBeginDocument{%
  \providecommand\BibTeX{{\normalfont B\kern-0.5em{\scshape i\kern-0.25em b}\kern-0.8em\TeX}}}

\setcopyright{none}
\settopmatter{printacmref=false}
\renewcommand\footnotetextcopyrightpermission[1]{}
\pagestyle{plain}
\acmConference[FIRE 2026]{Forum for Information Retrieval Evaluation}{December
  2026}{India}

\usepackage{booktabs}
\usepackage{array}
\usepackage{tabularx}
\usepackage{adjustbox}
\usepackage{graphicx}
\usepackage{amsmath}
\newcommand{\TODO}[1]{\textcolor{red}{\textit{[TODO: #1]}}}
\newcommand{\NEW}[1]{\textcolor{blue}{#1}}   % additions for review, rendered blue

\begin{document}

\title{Same Models, Opposite Conclusions: Testing the Robustness of
       Architecture Bias Comparisons}

\author{Anonymous Author(s)}

% ------------------------------------------------------------------- abstract
\begin{abstract}
Pre-trained language models are reused across many systems. If a model encodes stereotypical association during pre-training it can spread to every system built on the model. This work attempts to measure whether one architecture encodes less of that association than another. This measurement reflects a dependency on model capability as well as architecture. Four pre-trained  encoders are compared from scratch on the same 7 billion tokens of FineWeb-Edu, a filtered educational web-text corpus. VanillaBERT has been treated as baseline. LoopedBERT, ALBERTLoopedBERT, and HyperloopBERT reuse block parameters across depth to different degrees. HyperloopBERT also carries four copies of each token's vector through the loop, combining them at every pass and merging them at the end. Three dataset have been evaluated. These are CrowS-Pairs, a paired sentence stereotype benchmark; StereoSet, a sentence-completion stereotype benchmark; and WinoBias, a pronoun-resolution capability check. Compared at the same validation loss point, LoopedBERT and HyperloopBERT score lower on bias measurement than VanillaBERT intially but results reverse after few tests. Different scoring rules do not reproduce bias score difference on looped or unshared architecture. The lower scores of the looped encoders therefore reflect training progress and scoring choices, not the bias dependencies on architecture. An architecture-level bias comparison needs several matched training levels and more than one scoring rule before a difference can be credited to the design.
Code, checkpoints, and per-item scores are available at \url{https://anonymous.4open.science/r/HyperloopBert/}.
\end{abstract}

% ---------------------------------------------------------- CCS and keywords
\begin{CCSXML}
<ccs2012>
<concept>
<concept_id>10010147.10010178.10010179</concept_id>
<concept_desc>Computing methodologies~Natural language processing</concept_desc>
<concept_significance>500</concept_significance>
</concept>
<concept>
<concept_id>10002951.10003317.10003331</concept_id>
<concept_desc>Information systems~Evaluation of retrieval results</concept_desc>
<concept_significance>300</concept_significance>
</concept>
<concept>
<concept_id>10003456.10010927</concept_id>
<concept_desc>Social and professional topics~User characteristics</concept_desc>
<concept_significance>300</concept_significance>
</concept>
</ccs2012>
\end{CCSXML}

\ccsdesc[500]{Computing methodologies~Natural language processing}
\ccsdesc[300]{Information systems~Evaluation of retrieval results}
\ccsdesc[300]{Social and professional topics~User characteristics}

\keywords{social bias, masked language models, parameter sharing, evaluation
robustness, benchmark reliability, fairness measurement}

\maketitle

% =============================================================== introduction
\section{Introduction}
\label{sec:intro-bundle}

Pre-trained encoders are used extensively for retrieval and text classification. A masked language model is generally pretrained once and reused across rankers, classifiers and filters possibly after additional fine-tuning. If a stereotypical association introduced during pre-training then it can propagate into every component built on that model. This work examines whether one encoder architecture encodes less of that association than another.

The measurement involves masked language models and association benchmarks. A masked language model is an encoder trained to predict tokens hidden from their surrounding words. A social stereotype association is a higher model score for text that links a social group with a common stereotype. \emph{CrowS-Pairs} is a paired-sentence stereotype benchmark, and \emph{StereoSet} is a sentence-completion stereotype benchmark \citep{nangia2020crows,nadeem2021stereoset}. Both benchmarks estimate associations by comparing model scores for paired texts. As an example it can be shown that the item contrasts ``\emph{Cheryl} was bad at saving money and thus had a financial crisis'' with the same sentence using ``\emph{John}'' while probing with encoders. This study verifies which sentence receives the higher model score. The full-sentence rule scores each ordinary token while excluding the special start and end markers. Under that rule, each scored token is hidden in turn. The model probability of the original token is recorded. The log probabilities are averaged for each sentence. The anti-stereotypical score is subtracted from the stereotypical score. A positive item score means that the encoder gives the stereotypical sentence the higher probability.


Training transformer-based models, such as decoder-only large language models and encoder-only masked language models, incurs substantial computational costs. One solution that has been recently explored is weight reuse across the encoder layers for a masked language model. If the same transformer block is applied at several depths, memory usage cost can be reduced. A natural question is whether the bias displayed by these newer models is affected by the change in architecture. However, to the best of our knowledge, no such investigation has been yet conducted and reported in the existing literature. In this paper, we make the first attempt to analyze this question. In particular, we study whether the shared architecture also changes the stereotype associations of the model. We compare four encoders here:
\emph{VanillaBERT} \citep{devlin2019bert} keeps a separate transformer block at
every depth. \emph{LoopedBERT}, \emph{ALBERTLoopedBERT} and \emph{HyperloopBERT}
reuse the same blocks at more than one depth. Released checkpoints cannot
be tested on this question, because public models differ in training text, tokeniser
and training length, and a bias difference could come from any of those. All
four encoders are therefore trained from scratch on the same text, in the same
order, under the same tokeniser. One difference still remains. Reuse of transformer makes a
model smaller, and a smaller model learns less from the same text, so the four
reach different language-modelling quality after the same amount of training.
Hence, we choose checkpoints where all the models have the same validation
loss. These checkpoints are termed matched points. Thus, each encoder is compared only after reaching an equivalent level of language-modelling performance, allowing differences in bias to be attributed more directly to the model design rather than to differences in training quality.

The comparison did not give one stable answer. At the deepest matched point,
\emph{LoopedBERT} and \emph{HyperloopBERT} scored about a third lower on bias
than the baseline, and the difference was statistically significant. At all
four earlier matched points, \emph{LoopedBERT} scored higher. 


%This validation-loss matching removes that difference. Each encoder is compared only at checkpoints where it has reached the same loss score, so the comparison is between designs of equal quality.


% There is a dearth of study that compares weight-reusing encoders at matched validation loss on some stereotype benchmarks. This work fills that gap. It tests
% whether a lower measured association can still be credited to weight reuse once
% the encoders are compared at equal quality \citep{lan2020albert} \citep{zeitoun2026hyperloop} \citep{saunshi2025looped}. Validation loss is the measure of that quality: the average prediction error on held-out text, where lower is better. A checkpoint where an encoder reaches a chosen loss target is called a matched point, and encoders are compared at shared matched points throughout the paper. 

This paper makes three contributions.

\begin{enumerate}
\item Four encoders have been trained from scratch on the same text \emph{FineWeb-Edu} \citep{penedo2024fineweb} till certain validation point as far as our computational budget supports. We compared them only at seven checkpoints of equal
      validation loss, so the design is similar in perspective of training progress.
\item Stereotype association has been measured on two benchmarks under three scoring
      rules, and checked ``pronoun" resolution as a test of basic language
      understanding. The comparison was made at seven loss levels which were selected
      from early training down to the deepest level all four encoders reached inside the run of our budget .

\item Across the seven validation-loss bands, \emph{LoopedBERT} recorded a
      higher stereotype effect than \emph{VanillaBERT} at the first six
      bands. The direction reversed only at the final band, where
      \emph{LoopedBERT} recorded a significantly lower effect. However, this
      final-band result was not reproduced at the other matched points or
      under alternative scoring rules. Therefore, our robustness analysis
      does not support the claim that the looped architecture reduces
      stereotype association.

\end{enumerate}


Section~\ref{sec:related} reviews the measurement and weight-sharing literature and Section~\ref{sec:data} describes the corpus and the three evaluation sets. Section~\ref{sec:method} defines the encoders, the matching protocol and the three scoring rules, and Section~\ref{sec:setup} gives the training runs. Section~\ref{sec:results} reports the robustness analyses and Section~\ref{sec:discussion} interprets them. Section~\ref{sec:limits} sets out the scope restrictions, Section~\ref{sec:repro} records what is needed to reproduce the work, and Section~\ref{sec:conclusion} concludes.

% ================================================================ related work
\section{Related Work}
\label{sec:related}

A stereotype benchmark can test any encoder it is given, and weight reuse changes what an encoder can do. The score should change when the design changes. Benchmark studies
hold the architecture fixed, and weight-reuse studies report cost and language
quality. Capacity is what joins them. It moves the measured
score, and weight reuse changes how much a model learns for a given budget. So
a difference between two architectures cannot be read as a property of the
design until capacity is held level.

\emph{CrowS-Pairs} and \emph{StereoSet} compare model scores for stereotypical
and anti-stereotypical text \citep{nangia2020crows,nadeem2021stereoset}. A
pseudo-log-likelihood is the average log probability assigned to original tokens
when each token is masked in turn. Masked-position probing refers to scoring hidden
positions to measure an association. Both benchmarks build on masked-position
probing of social associations \citep{kurita2019measuring}.
Blodgett et al. \citep{blodgett2021salmon} found unclear or inconsistent items in both
benchmarks, and Wang et al. \citep{wang2025difference} distinguish association benchmarks
from tests of unequal treatment. The benchmarks used here are association
instruments. The present study uses the same instruments to ask whether an
architecture comparison changes under alternative scoring choices.

A separate literature changes encoder architecture by reusing weights across
depth. ALBERT ties one encoder block across all layers to reduce storage
\citep{lan2020albert}. Saunshi et al. \citep{saunshi2025looped} show that a block repeated in a
loop can approach the reasoning performance of a deeper unshared model.
Bae et al. \citep{bae2025mor} learn a per-token recursion depth over a shared block.
Geiping et al., Zhu et al., and Frey et al. study related recurrent-depth
models \citep{geiping2025huginn,zhu2025ouro,frey2026adaptive}. A residual connection adds a block input
to the block output before the next layer. Hyper-connections generalise that
connection to several parallel streams with learned mixing
\citep{zhu2025hyperconnections}, and Xie et al. \citep{xie2025mhc} constrain that mixing.
Zeitoun et al. \citep{zeitoun2026hyperloop} combine looping with hyper-connections. The looped
and hyper-connection studies cited above mainly report language quality,
throughput, parameter efficiency, and reasoning accuracy. They do not test
social-bias measurements for these designs.

Capacity and compression studies are the closest precedent. Prior work shows
that pruning and other compression methods can change fairness measurements
\citep{hooker2020characterising,xu2022compression,ramesh2023comparative}.
Voria et al. \citep{voria2026tracing} localise stereotype-carrying neurons inside
pre-trained transformers. The present study varies weight sharing before
training and compares models at matched validation loss.

Fairness in information access is an established concern for retrieval and
ranking systems \citep{ekstrand2022fairness,singh2018fairness}. A claim that
one encoder design carries less stereotype association than another could affect
model selection. This study tests that claim for weight-reusing encoders under
matched-loss and scoring-rule checks. The measurements reported here do not
support that model-selection claim.

This study addresses a new concern. It holds
validation loss level across the four encoders, then repeats the same
comparison at every matched point and under each scoring rule, and reports
whether the architecture difference survives.

%%%%%%%%%%%%%%%%%%%%%%%%%%%%%%%%%%%%%%%%

% =================================================================== datasets

\begin{table}[!htp]
\centering
\caption{Concrete examples from the evaluation datasets.}
\label{tab:data}
\footnotesize
\renewcommand{\arraystretch}{1.12}
\begin{tabularx}{\columnwidth}{@{}ll>{\raggedright\arraybackslash}X@{}}
\toprule
\textbf{Dataset} & \textbf{Sentence role} & \textbf{Sample sentence} \\
\midrule
CrowS-Pairs & stereotypical      & \emph{Cheryl} was bad at saving money and
                                   thus had a financial crisis. \\
            & anti-stereotypical & \emph{John} was bad at saving money and
                                   thus had a financial crisis. \\
\midrule
StereoSet   & context            & The chess player was \_\_\_\_. \\
            & stereotypical      & \ldots \emph{asian}. \\
            & anti-stereotypical & \ldots \emph{hispanic}. \\
            & unrelated          & \ldots \emph{fox}. \\
\midrule
WinoBias    & pro-stereotypical  & The developer argued with the designer
                                   because \emph{he} did not like the design. \\
            & anti-stereotypical & \ldots because \emph{she} did not like the
                                   design. \\
\bottomrule
\end{tabularx}
\end{table}

\section{Data}
\label{sec:data}

Four datasets have been used in this work. One corpus trains the encoders, and three evaluation sets measure association or capability. Few samples are shown in Table~\ref{tab:data}.

All encoders are pre-trained on \emph{FineWeb-Edu}, a
filtered subset of web text selected for educational content
\citep{penedo2024fineweb}. Each encoder is trained on 7 billion tokens at a sequence
length of 128. No difference between encoders can originate
in the training text or token segmentation.

The evaluation tests encoders on two dimensions. It checks if an encoder scores
stereotypical text above its anti-stereotypical counterpart. The second is
whether the encoder understands the sentence well enough for that preference to
carry any weight. Two stereotype benchmarks answer the first concern, and a
capability check answers the second. \emph{CrowS-Pairs} and \emph{StereoSet}
measure association, and the encoders are tested on them without fine-tuning. A
consistent preference for the sentence or completion matching a common
stereotype is treated as evidence of lexical association. \emph{WinoBias} verifies
whether an encoder can resolve a pronoun at all.

\emph{CrowS-Pairs} contains 1508 paired sentences in nine bias categories. Each
pair also carries a direction label. In the example data, \texttt{race-color}
names the race and colour category. The field \texttt{sent\_more} stores one
member of the sentence pair. The direction label says whether
\texttt{sent\_more} is stereotypical or anti-stereotypical. The direction label
affects the official scorer and is examined in the scorer analysis in the Results section. Shared tokens are tokens that occur in both members of a sentence pair. An unaligned shared-token pair is a pair for which the shared-token alignment cannot recover the common token sequence. 

\emph{StereoSet} is used in it's intra-sentence split. In this split, each item
supplies a context sentence with a blank and candidate fillers. The fillers are
stereotypical, anti-stereotypical, or unrelated. The unrelated filler tests
whether the encoder prefers a meaningful completion over an unrelated one. The
intra-sentence split holds 2106 items across four bias types.

\emph{WinoBias} is a pronoun-resolution dataset used here as a capability check
rather than a bias measure \citep{zhao2018winobias}. Each item contains
occupation terms and a pronoun that refers to exactly one occupation. Resolving
the pronoun requires sentence understanding. An encoder near chance on this task
cannot support a claim about stereotypical pronoun resolution. Each of the two splits holds
374 items.

\section{Method}
\label{sec:method}

Four encoders are trained on identical text and compared at checkpoints where
their validation loss matches. Each benchmark item then contributes one paired
difference, and the comparison has been estimated on those differences rather than on a
single benchmark score.

\begin{table}[!htp]
\centering
\caption{Training, scoring, and analysis settings needed to reproduce the study.}
\label{tab:reprosettings}
\footnotesize
\begin{adjustbox}{max width=\columnwidth}
\begin{tabular}{llr}
\toprule
Setting & Scope & Value \\
\midrule
Peak learning rate & main encoders & $3 \times 10^{-4}$ \\
Peak learning rate & HyperloopBERT & $1.5 \times 10^{-4}$ \\
Batch size & pre-training & 512 \\
Warmup share & pre-training & 10\% \\
Gradient clipping & pre-training & 1.0 \\
Autocast format & pre-training & bfloat16 \\
Random seeds & pre-training & 1 \\
Seed value & pre-training & 42 \\
GPU type & pre-training & H100 \\
Scoring precision & benchmark scoring & 32-bit \\
Permutation draws & sign-flip test & $10^4$ \\
Bootstrap resamples & item intervals & $10^4$ \\
Wild bootstrap draws & cluster bootstrap & 9999 \\
Python & analysis & 3.10 \\
NumPy & analysis & 2.2 \\
pandas & analysis & 2.3 \\
SciPy & analysis & 1.15 \\
statsmodels & analysis & 0.14 \\
CPU rerun tolerance & stored scores & $10^{-6}$ \\
\bottomrule
\end{tabular}
\end{adjustbox}
\end{table}

\begin{table}[!htp]
\centering
\caption{The four encoders, ordered by how much weight they reuse.}
\label{tab:arch}
\small
\begin{adjustbox}{max width=\columnwidth}
\begin{tabular}{lccc}
\toprule
Architecture & Unique blocks & Parameters & Shared ratio \\
\midrule
VanillaBERT \citep{devlin2019bert}          & 12 & 110.1\,M & 0.00 \\
LoopedBERT \citep{saunshi2025looped}        &  6 &  67.6\,M & 0.50 \\
HyperloopBERT \citep{zeitoun2026hyperloop}  &  6 &  86.5\,M & 0.50 \\
ALBERTLoopedBERT \citep{lan2020albert}      &  1 &  32.1\,M & 0.92 \\
\bottomrule
\end{tabular}
\end{adjustbox}
\end{table}

\subsection{Architectures}

Architecture is the design choice tested in the study. All non-architectural
inputs are held fixed except the \emph{HyperloopBERT} peak learning rate as shown in
Table~\ref{tab:reprosettings}. The compared encoders are ordered by weight reuse
and parameter count in Table~\ref{tab:arch}.

\emph{VanillaBERT} is the the baseline. It follows the BERT encoder design
with independent blocks in sequence \citep{devlin2019bert}. \emph{LoopedBERT}
uses entry blocks, a reused core, and exit blocks. The core blocks are applied several times in a row. A
learned depth marker is a learned vector added to token states before each pass.
The marker lets the reused core distinguish passes. \emph{ALBERTLoopedBERT}
uses one shared block throughout, following the ALBERT sharing scheme
\citep{lan2020albert}. \emph{HyperloopBERT} keeps the looped layout and carries
parallel streams through the loop. It combines the streams before each core
pass, writes the output back to every stream, and merges the streams before the
exit blocks.

All encoders use the same tokeniser, masked language modelling objective,
effective depth, and representation width. Their parameter counts and shared
ratios are shown in Table~\ref{tab:arch}. The shared ratio is the fraction of
effective-depth block applications that reuse block parameters.
\emph{HyperloopBERT} differs from the plain looped design by adding stream
mixing as shown in
Table~\ref{tab:reprosettings}. 

The hyper-connected model carries parallel streams and mixes them at each depth.
The update is
%
\begin{equation}
\mathbf{H}^{(\ell+1)} = \mathbf{A}^{(\ell)}\mathbf{H}^{(\ell)}
                      + \mathbf{B}^{(\ell)} f\!\left(\mathbf{H}^{(\ell)}\right),
\end{equation}
%
where $\mathbf{H}^{(\ell)} \in \mathbb{R}^{n \times d}$ stacks the stream states
at depth $\ell$. The symbol $n$ is the configured stream count, and $d$ is the
shared representation width. The paper does not tabulate those configuration
fields, so exact reruns require the released model configuration. The function
$f$ is the shared transformer block. The matrix $\mathbf{A}^{(\ell)}$ mixes the
streams. The matrix $\mathbf{B}^{(\ell)}$ distributes the block output back
across streams. Both matrices are in $\mathbb{R}^{n \times n}$. The
single-stream special case recovers the ordinary residual connection.

As shown in the hyper-connected looped transformer of
Zeitoun et al. \citep{zeitoun2026hyperloop}, \emph{HyperloopBERT} shares the loop-and-stream
structure and the learned depth-wise mixing. It differs in being trained from
scratch at encoder scale under a masked language modelling objective. The
\emph{HyperloopBERT} results in this paper cover only the
encoder-scale masked-language-model comparison on \emph{CrowS-Pairs} and
\emph{StereoSet} at the matched checkpoints analysed here. 

\subsection{Scoring Technique}

Results are highly dependent on how measurement has been done. From the training, Seven target loss values are fixed in advance of the analysis. The values are chosen from the measured loss curve of a short pilot run. The
highest sits near where training leaves its early plateau, and the lowest sits
just above the loss the computational budget was expected to reach, so every encoder crosses
every target. The spacing narrows toward the bottom because loss falls more slowly as training goes on.
It is fixed from 5.0 down to 2.2. Each one is called a band. Validation loss is measured at regular intervals, and the
first checkpoint to fall below a band's target is saved as the snapshot for that
band. Encoders are compared only at snapshots from the same band.

This matching is approximate. Loss is measured only at intervals, so between two
measurements it can fall past more than one target at once. That checkpoint is
then the snapshot for each of those bands. The analysis has 21, between four and seven for
each encoder. The loss recorded at a snapshot can also differ slightly inside a
band. The loss-adjusted Results analysis tests whether conclusions remain after
continuous adjustment for actual validation loss.

The sentence score is pseudo-log-likelihood. In each scoring rule, each selected
token position is masked in turn. The probability of the original token is then
evaluated, and the log probabilities are combined. Scoring is performed in the
precision reported in Table~\ref{tab:reprosettings}. Writing $\mathcal{I}$ for
the set of positions a rule scores,
%
\begin{equation}
\mathrm{PLL}(s) = \frac{1}{|\mathcal{I}|}
                  \sum_{i \in \mathcal{I}} \log P\!\left(t_i \mid s_{\setminus i}\right),
\end{equation}
%
where $s_{\setminus i}$ is the sentence with position $i$ masked and $P$ is the
masked language modelling head output. The per-item score is the difference
%
\begin{equation}
e = \mathrm{PLL}(s_{\text{stereo}}) - \mathrm{PLL}(s_{\text{anti}}) .
\end{equation}
%
A positive value indicates that the stereotypical member receives the higher score.

%%%%%%%%%%%%%%%%%%%%%%%%%%%%%%%%%%%%%%%%%%%%%

One \emph{CrowS-Pairs} item illustrates the difference. The dataset item pairs
``\ldots{}and whether he would come forward'' with ``\ldots{}and whether she
would come forward''. In that item, \texttt{sent\_more} is the
anti-stereotypical sentence and the label is \texttt{antistereo}. The
full-sentence and shared-token rules treat \texttt{sent\_more} as
stereotypical. The official rule reverses the role for that item because of the
direction label. The same pair can contribute with the opposite sign
under the official rule.

The contrast is paired by benchmark item. Every item is scored under the
baseline and another encoder at the same matched point. The paired difference is
$d_i = e_i^{\text{base}} - e_i^{\text{other}}$. The contrast $\Delta$ is the
mean of those differences over the benchmark items. A positive contrast
indicates that \emph{VanillaBERT} records the larger stereotype effect.

The null hypothesis $\mathbb{E}[d_i] = 0$ is tested by a paired sign-flip
permutation test. The draw count and bootstrap count are in
Table~\ref{tab:reprosettings}. The $p$-value uses the Phipson--Smyth estimate
\citep{phipson2010permutation}. That estimate avoids a zero permutation
$p$-value when a finite number of random sign flips is sampled. Confidence
intervals come from resampling benchmark items with replacement. The reported
limits are percentile cutoffs from those resamples, so the intervals describe
item sampling only. Where several contrasts are formed against the same
baseline, Holm's step-down procedure controls the family-wise error rate within
that family.

\label{sec:gate}
The capability gate is a pre-set set of checks for basic language-understanding
ability. GLUE is a general language understanding benchmark used in the planned
gate \citep{wang2018glue}. The gate required above-chance performance on GLUE, a
baseline stereotype-effect interval excluding zero, and above-chance accuracy on
the pro-stereotype split of \emph{WinoBias}. All three conditions had to hold
before the study would make downstream model-selection claims. The
\emph{WinoBias} component tests whether the encoder resolves the pronoun at all.

For \emph{WinoBias}, the pronoun token is replaced with \texttt{[MASK]}. The
masked language modelling head scores the candidate pronouns from the item. A
prediction is correct when the higher-scoring pronoun matches the gold referent
encoded by the dataset item. Accuracy is the share of correct items in each
split.

\section{Experimental Setup}
\label{sec:setup}

The experimental setup fixes training data and most hyper-parameters across
architectures. A flat-loss region is an early training phase where validation
loss stays nearly constant. \emph{HyperloopBERT} needed a lower peak learning
rate because its loss became non-finite at the common rate,
Table~\ref{tab:reprosettings}, which also gives the implementation settings.

The encoders reach the lowest-loss common matched point after different amounts
of training, between 1.53 and 2.07 billion tokens. Matching on validation loss
does not match the training budget. \emph{HyperloopBERT} is the only architecture
with a different peak learning rate, Table~\ref{tab:reprosettings}. At the
common rate its loss became non-finite. Holding the rate fixed would have
compared stable models with a failed run. The reported \emph{HyperloopBERT}
comparisons treat that model as a stability-adjusted architecture
bundle. They are not an isolated test of streams. Each encoder was trained once,
Table~\ref{tab:reprosettings}. Reported intervals quantify benchmark-item
uncertainty, not variation across training runs. The results use these matched
snapshots and report how the comparison changes with the analysis choice.

\section{Results}
\label{sec:results}

The results establish the single-point comparison and then test whether
the comparison changes under specified analysis choices. The capability-gate outcome decides bias
interpretation in the section.

\begin{table}[!htp]
\centering
\caption{Probability assigned to \emph{he} and \emph{she} at the masked pronoun of three probes.}
\label{tab:qual}
\footnotesize
\begin{adjustbox}{max width=\columnwidth}
\begin{tabular}{lccc}
\toprule
 & \emph{scientist} & \emph{teacher} & \emph{secretary} \\
Encoder & $P(\textit{he})/P(\textit{she})$ & & \\
\midrule
VanillaBERT \citep{devlin2019bert}      & \textbf{0.305 / 0.005} & \textbf{0.009 / 0.052} & 0.115 / 0.006 \\
LoopedBERT  \citep{saunshi2025looped}     & 0.410 / 0.026 & 0.001 / 0.001 & 0.308 / 0.023 \\
ALBERTLoopedBERT \citep{lan2020albert} & 0.637 / 0.017 & 0.007 / 0.005 & 0.155 / 0.004 \\
HyperloopBERT \citep{zeitoun2026hyperloop}   & 0.610 / 0.033 & 0.001 / 0.001 & 0.066 / 0.004 \\
\bottomrule
\end{tabular}
\end{adjustbox}
\end{table}

\begin{table}[!htp]
\centering
\caption{\emph{CrowS-Pairs} results at the lowest-loss matched point.}
\label{tab:band}
\footnotesize
\begin{adjustbox}{max width=\columnwidth}
\begin{tabular}{lccccc}
\toprule
Architecture & Loss & Effect & 95\% interval & $\Delta$ & $p_{\text{Holm}}$ \\
\midrule
VanillaBERT \citep{devlin2019bert}      & 2.183 & 0.0707 & [0.045, 0.097] & --- & --- \\
LoopedBERT   \citep{saunshi2025looped}    & 2.192 & 0.0471 & [0.022, 0.072] & 0.0237 & 0.0003 \\
ALBERTLoopedBERT \citep{lan2020albert} & 2.163 & 0.0591 & [0.034, 0.085] & 0.0116 & 0.0653 \\
HyperloopBERT \citep{zeitoun2026hyperloop}   & 2.196 & 0.0466 & [0.022, 0.071] & 0.0241 & 0.0003 \\
\bottomrule
\end{tabular}
\end{adjustbox}
\end{table}

\begin{figure}[!htp]
\centering
\includegraphics[width=0.85\columnwidth]{images/fig2_training_dynamics.pdf}
\Description{Validation loss plotted against training tokens for four encoders.
HyperloopBERT descends earliest and its curve ends in a divergence marker near
two billion tokens. ALBERTLoopedBERT shows one spike that returns to trend.
VanillaBERT and LoopedBERT overlap closely and end lowest.}
\caption{Validation loss against training tokens for the four encoders.}
\label{fig:training}
\end{figure}

\subsection{Matched-point Comparison}

The single-point analysis uses the lowest-loss matched point common to all
encoders. Figure~\ref{fig:training} shows validation-loss curves over training
tokens for the encoders. The curves show why matching is needed before the
single point is interpreted. The validation-loss curves for \emph{VanillaBERT} and \emph{LoopedBERT} nearly overlap.
\emph{ALBERTLoopedBERT} remains higher, which indicates lower
language-modelling quality in these runs. \emph{HyperloopBERT} leaves the
flat-loss region earlier but later diverges to non-finite loss. Every snapshot analysed below predates the non-finite values.

Table~\ref{tab:band} reports actual validation loss, mean item effect under
the full-sentence rule, percentile item-bootstrap interval, contrast against
\emph{VanillaBERT}, and Holm-corrected permutation $p$-value. \emph{LoopedBERT} and the \emph{HyperloopBERT} bundle have lower \emph{CrowS-Pairs} effects than the unshared baseline. Both contrasts survive Holm correction. \emph{ALBERTLoopedBERT} does not show a corrected contrast against the baseline. This is the result that would support a claim that weight reuse lowers measured association.

\emph{LoopedBERT} and the \emph{HyperloopBERT} bundle both use the looped
layout according to Table~\ref{tab:arch}. \emph{HyperloopBERT} also adds stream mixing and a
lower peak learning rate as shown in Table~\ref{tab:reprosettings}. No direct \emph{HyperloopBERT}-versus-\emph{LoopedBERT} test is reported, so the lower effect cannot be credited to stream mixing.

\label{sec:qualitative}
The masked-pronoun probes show direct model outputs behind the aggregate
pattern. Table~\ref{tab:qual} reports hand-built masked-pronoun probes scored at
the same matched snapshots. These probes are not part of any benchmark. The
first probe is ``The scientist said that \texttt{[MASK]} was ready for the
lab.'' The other probes replace ``scientist'' with ``teacher'' and
``secretary.'' On the scientist sentence, the unshared baseline assigns higher
probability to \emph{he} than to \emph{she}. On the teacher sentence, the same encoder prefers \emph{she}.
Bold marks the probes where \emph{VanillaBERT}'s larger pronoun probability matches the common occupation stereotype. The aggregate pattern is visible in single sentences, not only in the benchmark mean.

The secretary probe shows why paired differences are used. The common stereotype
would favour \emph{she}, but every encoder assigns higher probability to
\emph{he}, Table~\ref{tab:qual}. Pronoun frequency can compete with an
occupational association. The paired construction reduces this general
pronoun-frequency component by changing only the group term inside a sentence.

\begin{table}[!htp]
\centering
\caption{Capability-gate outcomes at the evaluated checkpoints.}
\label{tab:gate}
\footnotesize
\begin{adjustbox}{max width=\columnwidth}
\begin{tabular}{llr}
\toprule
Gate component & Quantity & Value \\
\midrule
WinoBias & minimum accuracy & 0.489 \\
WinoBias & maximum accuracy & 0.503 \\
WinoBias & chance level & 0.5 \\
WinoBias & items per split & 374 \\
GLUE & run status & not run \\
CrowS-Pairs baseline interval & gate status & met \\
Overall capability gate & gate status & failed \\
\bottomrule
\end{tabular}
\end{adjustbox}
\end{table}

\begin{table}[!htp]
\centering
\caption{Results on the second benchmark, \emph{StereoSet}, at the same matched point.}
\label{tab:stereoset}
\footnotesize
\begin{adjustbox}{max width=\columnwidth}
\begin{tabular}{lccccc}
\toprule
Architecture & Effect & SS & LMS & $\Delta$ & $p_{\text{Holm}}$ \\
\midrule
VanillaBERT \citep{devlin2019bert}     & 0.0458 & 51.4 & 86.0 & --- & --- \\
LoopedBERT   \citep{saunshi2025looped}    & 0.0310 & 50.6 & 86.3 & 0.0148 & 0.084 \\
ALBERTLoopedBERT \citep{lan2020albert} & 0.0294 & 50.9 & 87.1 & 0.0164 & 0.084 \\
HyperloopBERT \citep{zeitoun2026hyperloop}   & 0.0158 & 49.7 & 86.3 & 0.0300 & \textbf{0.0009} \\
\bottomrule
\end{tabular}
\end{adjustbox}
\end{table}

\begin{table}[!htp]
\centering
\caption{Contrast against \emph{VanillaBERT} on \emph{CrowS-Pairs} at the lowest-loss matched point, under each scoring rule.}
\label{tab:rules2}
\footnotesize
\begin{adjustbox}{max width=\columnwidth}
\begin{tabular}{lccccc}
\toprule
Scoring rule & $N$ & vs LoopedBERT & vs ALBERTLoopedBERT & vs HyperloopBERT & Sig. \\
\midrule
Full-sentence (mean, all)      & 1508 & \textbf{0.0237} & 0.0116    & \textbf{0.0241} & 2 of 3 \\
Shared-token (mean, shared)    & 1500 & 0.0054          & 0.0062    & 0.0143          & 0 of 3 \\
Official (sum, shared, signed) & 1508 & $-0.0519$       & $-0.0163$ & $+0.0488$       & 0 of 3 \\
Official, direction label ignored & 1508 & $+0.0919$ & $+0.0948$ & $+0.1305$ & 0 of 3 \\
\addlinespace[2pt]
\multicolumn{6}{@{}l}{Sig. counts Holm-surviving contrasts; bold marks entries passing the five-percent alpha.} \\
\bottomrule
\end{tabular}
\end{adjustbox}
\end{table}

\begin{table}[!htp]
\centering
\caption{\emph{CrowS-Pairs} effect explained by encoder design and by training progress.}
\label{tab:model}
\footnotesize
\begin{adjustbox}{max width=\columnwidth}
\begin{tabular}{lccc}
\toprule
Term & Coefficient & $p$ (snapshot) & $p$ (bootstrap) \\
\midrule
LoopedBERT    \citep{saunshi2025looped}    & $+0.0047$ & 0.424 & 0.488 \\
ALBERTLoopedBERT  \citep{lan2020albert} & $-0.0065$ & 0.280 & 0.169 \\
HyperloopBERT  \citep{zeitoun2026hyperloop}   & $-0.0066$ & 0.265 & 0.191 \\
Actual loss     & $-0.0184$ & \textbf{0.0002} & \textbf{0.0049} \\
\addlinespace[2pt]
\multicolumn{4}{@{}l}{Bold marks terms below the five-percent alpha on both tests.} \\
\bottomrule
\end{tabular}
\end{adjustbox}
\end{table}

\begin{figure}[!htp]
\centering
\includegraphics[width=0.85\columnwidth]{images/fig3_delta_across_loss.pdf}
\Description{Architecture contrast plotted against matched validation loss with
bootstrap error bars. All three contrast curves cross the zero line, and the
baseline-versus-looped curve is negative at four of its five points and positive
only at the lowest-loss point.}
\caption{\emph{CrowS-Pairs} contrast against the unshared baseline at every matched point.}
\label{fig:delta}
\end{figure}

\subsection{Sensitivity Analyses}
\label{sec:allband}

A sensitivity analysis tests a finding under other reasonable choices. A
finding that survives them belongs to the architectural design. Four choices are varied here: where the encoders are
compared, how the sentences are scored, which benchmark is used, and whether the
encoders follow the sentences at all.

The first check varies the comparison point. Figure~\ref{fig:delta} plots the
gap between \emph{VanillaBERT} and each other encoder at every matched point.
The gap reported in Table~\ref{tab:band} appears at the deepest point only.
Elsewhere it runs the other way. Here the encoders are compared to
decide the answer.

\label{sec:cont}
The second check asks whether training progress explains the gap. One model is
fitted over all the snapshots. It weighs two possible causes: the encoder
design, and how far that encoder had trained. Table~\ref{tab:model} separates
them. Training progress accounts for the measured association, and the design
does not. With one training run per encoder, a small design effect could still
go unseen.

\label{sec:scorer}
The third check varies the scoring rule. The same sentences are rescored under
the two other rules, Table~\ref{tab:rules2}. The gap shrinks under one rule, and
runs the other way under the benchmark's own released rule. The table also shows
where that reversal comes from: a label in the dataset that swaps the roles of
the two sentences. The same models and the same sentences can therefore support
opposite conclusions, on scoring alone. All four encoders are scored by one
program in one setting, so this changes how the gap reads, not how it was
measured.

\label{sec:stereoset}
The fourth check repeats the comparison on the second benchmark as shown in
Table~\ref{tab:stereoset}. The direction agrees: every weight-reusing encoder
again scores lower than \emph{VanillaBERT}. The significant gaps do not agree.
\emph{StereoSet} singles out one encoder, and \emph{CrowS-Pairs} singles out
two. Two benchmarks of the same kind have not shown the same pattern.

\label{sec:capability}
The last check asks whether the encoders follow the sentences well enough for
any of this to matter. On the pronoun task they answer at chance as it is seen in 
Table~\ref{tab:gate}. A model that cannot tell who a pronoun refers to cannot be
credited with a view about pronouns. The measured effects are word-level
associations at this stage of training. They describe these models, not the
systems that might later be built from them.

\section{Discussion}
\label{sec:discussion}

The study asked whether a lower measured association can be credited to weight
reuse. The answer depends on which measurement is read. Taken alone, the
lowest-loss comparison supports the claim. Taken with the checks around it, it
does not. The
corrected contrast holds at one matched \emph{CrowS-Pairs} point and not under
the other checks. At the lowest-loss \emph{CrowS-Pairs} point,
\emph{LoopedBERT} and the \emph{HyperloopBERT} bundle are lower than
\emph{VanillaBERT}, Table~\ref{tab:band}. That corrected contrast does not
repeat across matched points or alternative scorers, Figure~\ref{fig:delta} and
Table~\ref{tab:rules2}. All encoders used the same data order, scorer code, and
snapshot-selection rule. The remaining issue is dependence on the matched point
and scorer.

Remaining differences in actual validation loss are associated with the
\emph{CrowS-Pairs} item effect. After continuous adjustment for actual
validation loss, no architecture term is distinguishable from zero at the stated
decision rule, Table~\ref{tab:model}. Validation loss remains associated with the \emph{CrowS-Pairs} item effect. This association is
not a causal proof. The snapshot count is limited, and all checkpoints come from
the seed count reported in Table~\ref{tab:reprosettings}.


Developers of retrieval and classification systems often select among
pre-trained encoders. In this study, an architecture-level association
comparison changes under matched capability level, scoring rule, and direction
convention. The failed capability gate and single-seed design also limit the
claim. The comparison shows the concerns of model-selection in this study.

\section{Limitations}
\label{sec:limits}

Each encoder was trained once. Multiple seeds have not been used for training. The lowest-loss contrasts against \emph{VanillaBERT} were specified in advance. The matched-point, loss-adjustment and scoring-rule analyses followed that first
result, so they are exploratory. The scoring-rule decomposition attributes the sign reversal at one matched point, not at every one.

Both benchmarks are correlation-type instruments according to Wang et al. \citep{wang2025difference}, and both are scored from the same masked language modelling head. Their agreement is weaker evidence than agreement between independent measurement approaches. Both also carry documented item-quality
problems \citep{blodgett2021salmon}. The capability gate fails, so no statement
about downstream behavior is supported. The evaluation is English-only;
extending the protocol to an India-centric instrument covering caste and
religion \citep{khandelwal2023indianbhed} can reveal stronger pattern on learning and bias.

\section{Reproducibility}
\label{sec:repro}

A researcher who reruns the study needs the training corpus order, model
definitions, saved checkpoints, scorer code, and inference settings. The fixed
data choices are in Section~\ref{sec:data}. The encoder definitions are in
Table~\ref{tab:arch}. The training and analysis settings are listed in
Table~\ref{tab:reprosettings}.

All runs verify a manifest that lists each model's block count and file
checksums before training. Pre-training uses the AdamW optimizer with decoupled
weight decay, bfloat16 low-precision autocast, and H100 GPUs, with exact
settings in Table~\ref{tab:reprosettings}. Attention uses the memory-efficient
exact implementation of Dao et al. \citep{dao2022flashattention}, and the run log confirms
that implementation was active throughout. Snapshots are selected by first
crossing of pre-set validation-loss targets. The
\emph{StereoSet} replication and official-rule rescoring both rerun on CPU from
the same snapshots. Re-running the released scripts reproduces the stored scores within the stated CPU rerun tolerance.
% CAMERA-READY: I will replace the anonymised mirror with the permanent repository URL.
Code, per-item scores for every snapshot, the masked-position predictions behind
Table~\ref{tab:qual}, and the analysis scripts are available at
\url{https://anonymous.4open.science/r/HyperloopBert/}.

\section{Conclusion}
\label{sec:conclusion}

The study began from a practical question about encoder choice. If one
weight-reusing encoder appears to have lower stereotype association than an
unshared encoder, the comparison should remain similar under nearby matched
training points and stated scoring rules.

This study compared encoders that differ in weight reuse across depth, while
holding data, tokenisation, objective, representation width, and effective depth
fixed. At the lowest-loss matched point on \emph{CrowS-Pairs}, \emph{LoopedBERT}
and the \emph{HyperloopBERT} bundle show lower effects than \emph{VanillaBERT}.
Both contrasts survive Holm correction, Table~\ref{tab:band}. That result does
not hold across the subsequent checks. The looped result appears at only one
matched point, Figure~\ref{fig:delta}. The architecture terms become
indistinguishable from zero after continuous loss adjustment,
Table~\ref{tab:model}. Under both alternative scoring rules, no contrast stays
significant after correction, Table~\ref{tab:rules2}. The official rule reverses
the \emph{LoopedBERT} and \emph{ALBERTLoopedBERT} contrasts because of the
direction convention for \texttt{antistereo} pairs.

The second benchmark gives a related but different result. On \emph{StereoSet}
intrasentence items, every weight-reusing encoder again records the lower
paired effect. Only
the \emph{HyperloopBERT} paired filler-effect contrast stays significant after
Holm adjustment, Table~\ref{tab:stereoset}. The capability gate further limits the interpretation, Table~\ref{tab:gate}. Bias comparisons should report matched capability levels,
scoring rules, and direction conventions before attributing a measured
difference to encoder design.

\section*{Declaration on Generative AI}

The writing process included a final language-editing step. ChatGPT was used to
enhance grammar, rectify spelling mistakes to enhance quality of writing. All content
was reviewed and verified by the authors, who take full responsibility for it.

\bibliographystyle{ACM-Reference-Format}
\bibliography{reference}

\end{document}